\chapter{Sicherheit}
\label{ch:quality/security}

Sicherheit ist sehr wichtig, damit dein Mod nicht dazu führt, dass der
Serverbetreiber Daten oder die Kontrolle verliert.

\section{Grundprinzipien}\label{grundprinzipien}

Das wichtigste Prinzip der Sicherheit lautet: \textbf{Vertraue niemals
dem Benutzer.} Alles, was ein Benutzer übermittelt, ist als bösartig zu
behandeln. Prüfe also immer, ob seine Eingaben gültig sind, ob er die
nötigen Rechte hat und ob er die Aktion auch sonst ausführen darf (also
etwa in Reichweite oder Besitzer ist).

Eine bösartige Aktion muss nicht unbedingt Daten verändern oder
zerstören; sie kann auch auf sensible Daten zugreifen, etwa auf
Passwort-Hashes oder private Nachrichten. Besonders schlimm ist das,
wenn der Server Angaben wie E-Mail-Adressen oder das Alter speichert,
wie es manche zur Verifizierung tun.

\section{Formspecs}\label{formspecs}

\subsection{Eingaben niemals
vertrauen}\label{eingaben-niemals-vertrauen}

Jeder Benutzer kann fast jede Formspec\index{formspec@Formspec} mit
beliebigen Werten jederzeit abschicken.

Hier echter Code aus einem Mod:

\begin{Shaded}
\begin{Highlighting}[]
\VariableTok{core}\OperatorTok{.}\NormalTok{register\_on\_player\_receive\_fields}\OperatorTok{(}
    \KeywordTok{function}\OperatorTok{(}\VariableTok{player}\OperatorTok{,} \VariableTok{formname}\OperatorTok{,} \VariableTok{fields}\OperatorTok{)}
        \ControlFlowTok{for} \VariableTok{key}\OperatorTok{,} \VariableTok{field} \KeywordTok{in} \FunctionTok{pairs}\OperatorTok{(}\VariableTok{fields}\OperatorTok{)} \ControlFlowTok{do}
            \KeywordTok{local} \VariableTok{x}\OperatorTok{,} \VariableTok{y}\OperatorTok{,} \VariableTok{z} \OperatorTok{=} \FunctionTok{string.match}\OperatorTok{(}
                \VariableTok{key}\OperatorTok{,}
                \StringTok{"goto\_([\%d{-}]+)\_([\%d{-}]+)\_([\%d{-}]+)"}
            \OperatorTok{)}
            \ControlFlowTok{if} \VariableTok{x} \KeywordTok{and} \VariableTok{y} \KeywordTok{and} \VariableTok{z} \ControlFlowTok{then}
                \VariableTok{player}\OperatorTok{:}\NormalTok{set\_pos}\OperatorTok{(\{}
                    \VariableTok{x} \OperatorTok{=} \FunctionTok{tonumber}\OperatorTok{(}\VariableTok{x}\OperatorTok{),}
                    \VariableTok{y} \OperatorTok{=} \FunctionTok{tonumber}\OperatorTok{(}\VariableTok{y}\OperatorTok{),}
                    \VariableTok{z} \OperatorTok{=} \FunctionTok{tonumber}\OperatorTok{(}\VariableTok{z}\OperatorTok{),}
                \OperatorTok{\})}
                \ControlFlowTok{return} \KeywordTok{true}
            \ControlFlowTok{end}
        \ControlFlowTok{end}
    \KeywordTok{end}
\OperatorTok{)}
\end{Highlighting}
\end{Shaded}

Siehst du das Problem? Ein bösartiger Benutzer kann eine Formspec mit
eigenen Positionswerten abschicken und sich so an jeden beliebigen Ort
teleportieren. Mit modifizierten Clients ließe sich das sogar
automatisieren -- praktisch ein \texttt{/\allowbreak{}teleport} ganz
ohne Privileg.\index{privileg@Privileg}

Die Lösung für solche Probleme ist ein
Kontext~(Kapitel~\ref{ch:players/formspecs}), wie im Kapitel über
Formspecs gezeigt.

\subsection{Zeitpunkt der Prüfung ist nicht Zeitpunkt der
Verwendung}\label{zeitpunkt-der-pruxfcfung-ist-nicht-zeitpunkt-der-verwendung}

Jeder Benutzer kann jede Formspec mit beliebigen Werten jederzeit
abschicken, außer wo die Engine es verhindert:

\begin{itemize}
\tightlist
\item
  Die Eingabe einer Node-Formspec wird blockiert, wenn der Benutzer zu
  weit entfernt ist.
\item
  Seit Version 5.0 werden benannte Formspecs blockiert, die noch nicht
  angezeigt wurden.
\end{itemize}

Prüfe deshalb im Handler, ob der Benutzer die Bedingungen erfüllt, unter
denen die Formspec überhaupt angezeigt wird, und ebenso die Bedingungen
für die jeweiligen Aktionen.

Die Sicherheitslücke, die entsteht, wenn Rechte beim Anzeigen der
Formspec, aber nicht beim Verarbeiten geprüft werden, heißt \emph{Time
of Check is not Time of Use} (TOCTOU).

\section{(Unsichere) Umgebungen}\label{unsichere-umgebungen}

Luanti erlaubt Mods, eine Umgebung ohne Sandbox anzufordern, die Zugriff
auf die vollständige Lua-API gibt.

Siehst du die Sicherheitslücke im folgenden Code?

\begin{Shaded}
\begin{Highlighting}[]
\KeywordTok{local} \VariableTok{ie} \OperatorTok{=} \VariableTok{core}\OperatorTok{.}\NormalTok{request\_insecure\_environment}\OperatorTok{()}
\VariableTok{ie}\OperatorTok{.}\VariableTok{os}\OperatorTok{.}\NormalTok{execute}\OperatorTok{((}\StringTok{"path/to/prog \%d"}\OperatorTok{):}\FunctionTok{format}\OperatorTok{(}\DecValTok{3}\OperatorTok{))}
\end{Highlighting}
\end{Shaded}

\texttt{string.\allowbreak{}format} ist eine Funktion in der global
geteilten Tabelle \texttt{string}. Ein bösartiger Mod könnte diese
Funktion überschreiben und so etwas an \texttt{os.\allowbreak{}execute}
übergeben:

\begin{Shaded}
\begin{Highlighting}[]
\FunctionTok{string.format} \OperatorTok{=} \KeywordTok{function}\OperatorTok{()}
    \ControlFlowTok{return} \StringTok{"xdg{-}open \textquotesingle{}http://example.com\textquotesingle{}"}
\KeywordTok{end}
\end{Highlighting}
\end{Shaded}

Der Mod könnte dabei viel Bösartigeres übergeben, als eine Website zu
öffnen, etwa einem entfernten Angreifer die Kontrolle über den Rechner
geben.

Einige Regeln für unsichere Umgebungen:

\begin{itemize}
\tightlist
\item
  Speichere sie immer in einer lokalen Variablen und übergib sie nie an
  eine Funktion.
\item
  Stelle sicher, dass du allen Eingaben an unsichere Funktionen
  vertrauen kannst, um das obige Problem zu vermeiden. Vermeide also
  global überschreibbare Funktionen.
\end{itemize}
